% Appendix to Section 3 (learning one schema).  Source: research/tracks/single/notes-final.md (+ referee.md).
% Uses the local macros \single... defined at the top of sections/single.tex.

\section{Proofs and further results for Section~\ref{sec:single}}\label{app:single}

This appendix gives the full proofs of \Cref{sec:single}, in the notation of \Cref{sec:single:prelim}, and the secondary results to which the main text only points. Bound variables are de Bruijn indices. A datum $d$ in $\inst(T)$ has the unique matcher $\theta_d$, and $\beta^d_M=\theta_d(M)$. All proofs were transcribed from the research record and checked again; we found no error, and some steps are spelled out in more detail than there.

% =====================================================================
\subsection{Basic lemmas}\label{app:single:basic}

Parts (a), (b) and (e) of \Cref{lem:single:basic} are \Cref{lem:setting:preservation,lem:setting:subst,lem:setting:renaming}, proved in \Cref{app:setting:lemmas,app:setting:renaming}.

\begin{proof}[Proof of \Cref{lem:single:basic}(c), (d)]
(c) The skeleton is prefix-closed: the parent of a skeleton position is a skeleton position, because the children of an occurrence are its arguments. By \Cref{lem:setting:preservation}, every datum carries the symbol $T|_p$ at each $p\in\singleskel(T)$; by induction on the depth of $p$, $p\in C(D)$. The parent of an occurrence is in $\singleskel(T)\subseteq C(D)$, so the occurrence lies in $C(D)$ or is a slot.

(d) Match $d|_\sigma$ against $d|_r$, treating the free indices of $d|_\sigma$ as first-order variables. Walk both in lock step. At a free-index occurrence $y$ of $d|_\sigma$ under $e$ local binders, read off the subterm of $d|_r$ at the same place and lower it by $e$; fail if it mentions one of those $e$ local binders, or if two occurrences of $y$ disagree; fail on any other symbol clash. The walk either fails or returns a partial map $e_d$ defined exactly on $\FV(d|_\sigma)$, and $d|_r=(d|_\sigma)[\bar u]$ holds iff $\bar u$ extends $e_d$. So a common $\bar u$ on $Y$ exists iff no $e_d$ fails and the $e_d$ agree on common arguments. This is a pairwise condition, and $\bar u$ is then determined on $\bigcup_d\FV(d|_\sigma)=Y$. The walks take time $O(\sum_d|d|)$. The argument does not use finiteness of $D$.
\end{proof}

% =====================================================================
\subsection{Minimal covering templates}\label{app:single:min}

\begin{proof}[Proof of \Cref{thm:single:min}]
(a) We use two reduction steps. Each turns a covering $T\in\DT$ into a covering $T'=T\sigma\in\DT$, so $T\gen T'$.

\emph{Step V (vacuous place).} Suppose $z_m$ occurs in no $\beta^d_M$. Put
\[
\sigma(M):=\lambda\bar z.M'(z_1,\dots,z_{m-1},z_{m+1},\dots,z_n)\qquad(M'\text{ fresh}).
\] A pattern occurrence becomes an occurrence with a sub-list of distinct indices, again a pattern occurrence. $D$ is covered with $\theta'_d(M'):=\beta^d_M$ (holes renumbered).

\emph{Step H (common root).} Suppose all $\beta^d_M$ ($d\in D$) have the same root.
\begin{itemize}
\item If the root is a hole $z_m$ ($M$ is term-valued), put $\sigma(M):=\lambda\bar z.z_m$: each occurrence $M(\bar t)$ becomes $t_m$. $D$ is still covered, since $d|_q=z_m[\bar t]=t_m$ at each occurrence $q$.
\item If the root is a symbol $f$ with children $1,\dots,k$, child $j$ under $e_j\in\{0,1\}$ binders of $f$, put $\sigma(M):=\lambda\bar z.f(M_1(\bar z,0^{e_1}),\dots,M_k(\bar z,0^{e_k}))$, with fresh $M_j$ of arity $n+e_j$; here $0^{e}$ is the index $0$ (the variable bound by $f$) if $e=1$ and nothing if $e=0$, and the holes are shifted under $f$ by plugging. A pattern occurrence $M(\bar y)$ becomes $f(\dots,M_j(\bar y{\uparrow},0),\dots)$, whose arguments are distinct indices: a pattern occurrence. Other occurrences get metavariable-free arguments. $D$ is covered with $\theta'_d(M_j):=$ the $j$-th child of $\beta^d_M$, with the variable bound by $f$ replaced by the new hole. (If $f$ is a constant or a parameter, $k=0$ and $M$ disappears.)
\end{itemize}
Let $\mu(T)=(|C(D)|-|\singleskel(T)|,\ \sum_M\mathrm{arity}(M))\in\N^2$, ordered lexicographically; the first component is $\ge0$ by \Cref{lem:single:basic}(c). Step H strictly increases $|\singleskel|$: the root symbol, or the index $y_m$, becomes rigid at least at the pattern position. Step V keeps the skeleton and lowers $\sum_M\mathrm{arity}(M)$. So the process stops, and it stops exactly when (S1) and (S2) hold, at a saturated template. Generality is transitive.

(b) $\singleskel\subseteq C(D)$ is \Cref{lem:single:basic}(c). Let $M(\bar y)$ be a pattern occurrence at $\sigma$. By \Cref{lem:setting:preservation}, $d|_\sigma=\beta^d_M[\bar y]$; its root is the root of $\beta^d_M$ if that is a symbol, and $y_m$ if $\beta^d_M=z_m$ (a body has no free index, so an index root comes from a hole). Since the $y_m$ are distinct, the roots of the $d|_\sigma$ are all equal iff the roots of the $\beta^d_M$ are. By (S2) they are not, so $\sigma\notin C(D)$; its parent lies in $\singleskel\subseteq C(D)$, so $\sigma$ is a slot. By (S1) each $y_m$ is free in some $d|_\sigma$, so $\{\bar y\}\subseteq Y_\sigma(D)$; coverage forces $\FV(d|_\sigma)\subseteq\{\bar y\}$, so $Y_\sigma(D)=\{\bar y\}$. Every occurrence lies in $A(D)$ by \Cref{lem:single:basic}(c). For another occurrence $M(\bar t)$ at $q$, the value $\beta^d_M=(d|_\sigma)[\bar z/\bar y]$ is determined by $D$ and the pattern slot $\sigma$; by (S1) each hole $z_m$ occurs in some $\beta^d_M$, and by \Cref{lem:setting:subst}(ii), $t_m$ is determined by $d|_q=\beta^d_M[\bar t]$.

So, up to renaming metavariables and permuting argument places (\Cref{lem:setting:renaming}), a saturated covering template is determined by its set $Q\subseteq A(D)$ of occurrence positions and by the map that sends each $q\in Q$ to the pattern slot of the metavariable occurring there (one pattern occurrence fixed per metavariable). The skeleton is then determined too: it consists of the positions of $C(D)$ that are neither in $Q$ nor below an element of $Q$. The sort of each metavariable is the sort of its positions, and its arity is $|Y_\sigma(D)|$. This gives the bound $2^{|A(D)|}\cdot|\singleslots(D)|^{|A(D)|}$.

(c) Let $T$ be minimal. By (a), $T\gen T'$ with $T'$ saturated; minimality gives $T'\gen T$, and by \Cref{lem:setting:renaming} $T$ is a renaming of $T'$, hence saturated. Now let $T$ cover $D$ arbitrarily, and let $S_T$ be the set of saturated covering templates $T'$ with $T\gen T'$. It is finite up to renaming by (b) and nonempty by (a), and $\gen$ is a partial order on renaming classes (\Cref{lem:setting:renaming}). Take $T_m$ $\gen$-minimal in $S_T$. If $U$ covers $D$ and $T_m\gen U$, then $U\gen U'$ for some saturated $U'$, by (a), and $T\gen T_m\gen U\gen U'$; so $U'\in S_T$ and $T_m\gen U'$, and minimality in $S_T$ gives $U'\equiv T_m$, hence $U\equiv T_m$. So $T_m$ is minimal among all covering templates, and $T\gen T_m$.

(d) is (c) with $T=T^*$.
\end{proof}

\begin{example}[Induction data, no anchor]\label{ex:single:c82}\status{computed}\src{single \S4}
For $D=\{\Ind(x=x),\Ind(0=x)\}$ (both motives have root $=$, so $D$ is no anchor), $\Min(D)$ consists of exactly the four templates
\[
\bigl(\zeta_1=\zeta_2\wedge\forall x(f(x)=x\to f(Sx)=Sx)\bigr)\to\forall x\,f(x)=x,\qquad\zeta_1,\zeta_2\in\{0,f(0)\},
\]
of which only $\zeta_1=f(0)$, $\zeta_2=0$ lies below $T_{\Ind}$. The earlier analysis reported eight\src{prior induction C8.2}, from a search bounded by size $23$; the template with $\zeta_1=\zeta_2=f(0)$ has size $24$, and five of the eight lie strictly above it. Brute force at size $\le24$ and the referee's independent enumerator confirm the four.
\end{example}

\noindent\emph{In detail.} The $D$-features are: the Sym features of the common frame (including the whole antecedent $0=0$ and the right-hand sides $x$, $Sx$, $x$ of the three atoms in the motive copies under $\forall x$); Scope with scope $\{x\}$ at the three slots; and the Eq features $\sigma_b\leftrightarrow\sigma_d$ (identity), $\sigma_b\to\sigma_c$ and $\sigma_d\to\sigma_c$ ($x\mapsto Sx$), and $\sigma_b,\sigma_d\to$ each of the two zeros of the antecedent ($x\mapsto0$). (No feature leaves $\sigma_c$: its column $(Sx,0)$ cannot be mapped onto $(x,0)$, nor onto a rigid $0$ or $Sx$.) A minimal template therefore merges $\sigma_b$ and $\sigma_d$ into $f(x)$, puts $f(Sx)$ at $\sigma_c$, and chooses at each zero of the antecedent between the rigid $0$ and $f(0)$: four incomparable templates, separated by the instance $f:=\lambda z.Sz$. Computed: the saturated templates number $35$, of which $4$ are minimal; brute force at size $\le24$ finds the same $4$, and every one of $2749$ enumerated covering templates lies above one of them; at size $\le23$ there are $8$ apparent minima, as reported earlier.\src{single e2; referee t1, t1b}

% =====================================================================
\subsection{The feature verifier}\label{app:single:feat}

\begin{proof}[Proof of \Cref{thm:single:feat}]
\emph{The templates.} $T_0(D)$ and every $T_\varphi$ are in $\DT$: every $M_\sigma$ has the pattern occurrence $M_\sigma(Y_\sigma)$, and in $T_\varphi$ the metavariable $M_{\sigma_0}$ keeps its pattern occurrence at $\sigma_0$, which is not below $r$ and differs from $r$ (as $\sigma_0\perp r$). Both cover $D$ with $\theta_d(M_\sigma):=(d|_\sigma)[\bar z/Y_\sigma]$, which is closed except for holes because $\FV(d|_\sigma)\subseteq Y_\sigma(D)$; at $r$, $\theta_d(M_{\sigma_0})[\bar u]=(d|_{\sigma_0})[\bar u]=d|_r$ because $\varphi$ is a $D$-feature.

\emph{(1) $\Acc(D)\subseteq\Feat(D)=\inst(T_0(D))\cap\bigcap_\varphi\inst(T_\varphi)$.} A sentence lies in $\inst(T_0(D))$ iff it has every Sym and every Scope feature: such a $q$ has the skeleton $C(D)$, and $q=T_0(D)\theta$ with $\theta(M_\sigma):=(q|_\sigma)[\bar z/Y_\sigma]$; conversely an instance has the skeleton, and $\FV(\theta(M_\sigma)[Y_\sigma])\subseteq Y_\sigma$. If $q$ has every $D$-feature, then it has the skeleton of $T_\varphi$, Scope at every slot, and $q|_r=(q|_{\sigma_0})[\bar u]$, so $q=T_\varphi\theta$ with the same $\theta$. Conversely every instance $q=T_\varphi\theta$ has $\varphi$: $q|_{\sigma_0}=\theta(M_{\sigma_0})[Y_{\sigma_0}]$ and $q|_r=\theta(M_{\sigma_0})[\bar u]=(q|_{\sigma_0})[\bar u]$ by \Cref{lem:setting:subst}(i). So $\Feat(D)$ equals the displayed intersection, and $\Acc(D)\subseteq\Feat(D)$ because $T_0(D)$ and the $T_\varphi$ cover $D$.

\emph{(2) $\Feat(D)\subseteq\Acc(D)$.} Let $T$ cover $D$. By \Cref{thm:single:min}(a), $T\gen T'$ with $T'$ saturated and covering, so $\inst(T)\supseteq\inst(T')$. Fix for every metavariable $M$ of $T'$ a pattern occurrence $M(\bar y)$, at the slot $\sigma_M$ (\Cref{thm:single:min}(b)). A sentence $q$ lies in $\inst(T')$ if it has
(i) $\singleSym(p)$ for every $p\in\singleskel(T')$,
(ii) $\singleScope(\sigma_M)$ for every $M$, and
(iii) $\singleEq(\sigma_M,r,\bar u_r)$, with $\bar u_r:=(y_m\mapsto t_{r,m})_m$, for every other occurrence $M(\bar t_r)$ at $r$.
Indeed, by (ii) $\theta(M):=(q|_{\sigma_M})[\bar z/\bar y]$ is a body, and (i), (iii) give $T'\theta=q$, since $\theta(M)[\bar t_r]=(q|_{\sigma_M})[\bar u_r]=q|_r$. Each of these is a $D$-feature: (i) by \Cref{lem:single:basic}(c); (ii) because $\sigma_M$ is a slot with $Y_{\sigma_M}(D)=\{\bar y\}$; (iii) because $r\in A(D)$ (\Cref{lem:single:basic}(c)), $r\perp\sigma_M$ (distinct occurrences are incomparable), the sorts agree, $\bar u_r$ maps $Y_{\sigma_M}(D)$ to metavariable-free terms in scope at $r$, and $T'$ covers $D$: $d|_r=\beta^d_M[\bar t_r]=(\beta^d_M[\bar y])[\bar u_r]=(d|_{\sigma_M})[\bar u_r]$. So $\Feat(D)\subseteq\inst(T')\subseteq\inst(T)$ for every covering $T$, i.e.\ $\Feat(D)\subseteq\Acc(D)$.
\end{proof}

\paragraph{Relation to anti-unification.} $T_0(D)$ is the higher-order pattern anti-unifier without merging; the pattern lgg's merging of columns equal up to renaming \citep{pfenning1991unification,baumgartner2017higher} is the case of an Eq feature whose $\bar u$ is a permutation. $\DT$ adds \emph{derived} occurrences, and the price is that the version space can have exponentially many minimal elements (\Cref{prop:single:blowup}); their intersection is still described by polynomially many equations (\Cref{cor:single:poly}). Second-order generalization beyond patterns is also studied by \citet{hirata2004generalization} and in the frameworks surveyed by \citet{cerna2023antiunification}; we have not checked whether $\DT$ or \Cref{thm:single:feat} appears there.

\begin{proof}[Proof of \Cref{cor:single:poly}]
Positions in $A(D)$ are positions of every datum (a slot's parent carries the same symbol in all data, hence has the same arity), so $m=|A(D)|\le\min_d|d|$. $C(D)$, the slots and the scopes are computed in one simultaneous walk, in time $O(n)$. For each of the at most $m^2$ ordered pairs $(\sigma,r)$ with $\sigma$ a slot, $r\in A(D)$, $\sigma\perp r$, of the same sort, \Cref{lem:single:basic}(d) decides in time $O(n)$ whether an Eq-feature exists and returns its unique $\bar u$. So there are at most $m^2$ Eq-features, plus at most $m$ Sym and $m$ Scope features. Each feature is checked on $q$ in time $O(|q|)$: Sym and Scope directly; $\singleEq(\sigma,r,\bar u)$ by walking $q|_\sigma$ and $q|_r$ in lock step, comparing each free-index occurrence $y$ of $q|_\sigma$ with the corresponding subterm of $q|_r$ against $u_y$ (shifted), which consumes that subterm. Hence $O(m^2|q|)$ per query.
\end{proof}

% =====================================================================
\subsection{Anchors}\label{app:single:anchor}

\begin{proof}[Proof of \Cref{thm:single:anchor}]
($\Leftarrow$; no assumption.) By \Cref{lem:single:basic}(c), $\singleskel(T^*)\subseteq C(D)$, and by \evRs\ no occurrence lies in $C(D)$; since $C(D)$ is prefix-closed and the data positions are the skeleton, the occurrences and the positions below occurrences, $C(D)=\singleskel(T^*)$, $\singleslots(D)=\Occ(T^*)$ and $A(D)=\singleskel(T^*)\cup\Occ(T^*)$. For $\sigma$ carrying $M(\bar t_\sigma)$,
\[
Y_\sigma(D)=\bigcup_i\FV(\beta^i_M[\bar t_\sigma])=\bigcup_i\ \bigcup_{m:\ z_m\in\beta^i_M}\FV(t_{\sigma,m})=Y^*_\sigma
\]
by \evN. Let $q=T^*\theta$. Then $q$ has every Sym feature (\Cref{lem:setting:preservation}) and every Scope feature ($\FV(q|_\sigma)\subseteq Y^*_\sigma$). Let $\singleEq(\sigma,r,\bar u)$ be a $D$-feature; $(\sigma,r)$ is one of the pairs of \Cref{def:single:events}. By \evU\ it is valid: $\bar t_r=\bar t_\sigma[\bar u_0]$. Then $s_i|_r=\beta^i_M[\bar t_\sigma[\bar u_0]]=(s_i|_\sigma)[\bar u_0]$ for all $i$ (\Cref{lem:setting:subst}(i)), and by forcing (\Cref{lem:single:basic}(d)) $\bar u=\bar u_0$ on $Y_\sigma(D)=Y^*_\sigma$. So $q|_r=\theta(M)[\bar t_\sigma[\bar u]]=(\theta(M)[\bar t_\sigma])[\bar u]=(q|_\sigma)[\bar u]$. Hence $q\in\Feat(D)=\Acc(D)$ (\Cref{thm:single:feat}): $\inst(T^*)\subseteq\Acc(D)$, i.e.\ $D$ is an anchor.

($\Rightarrow$; under \Rich.) If $D$ is an anchor, every instance of $T^*$ has every $D$-feature (\Cref{thm:single:feat}). We exhibit an instance violating a $D$-feature whenever an event fails; metavariables not mentioned get arbitrary values.
\begin{itemize}
\item \evRs\ fails at $\sigma\in\Occ(T^*)$, carrying $M$, with common root $h$. The ancestors of $\sigma$ lie in $\singleskel(T^*)\subseteq C(D)$, so $\sigma\in C(D)$ and $\singleSym(\sigma)$ is a $D$-feature. By \Rich\ choose $\theta(M)$ with a symbol root $\neq h$ (a closed term, resp.\ a formula body); by \Cref{lem:setting:subst}(iii), $(T^*\theta)|_\sigma$ has that root.
\item \evRs\ holds and \evN\ fails at $(M,m)$. Take a pattern occurrence $M(\bar y)$ at $\sigma$; $\sigma$ is a slot. Since the $y$'s are distinct and the $\beta^i_M$ are closed except for holes, $y_m$ is free in $s_i|_\sigma=\beta^i_M[\bar y]$ only if $z_m$ occurs in $\beta^i_M$; so $y_m\notin Y_\sigma(D)$, and $\singleScope(\sigma)$ is a $D$-feature. Choose $\theta(M)$ containing $z_m$: $z_m$ itself, or $R(z_m,\dots,z_m)$ for a relation symbol $R$, which exists by \Rich. Then $y_m\in\FV((T^*\theta)|_\sigma)$.
\item \evRs\ and \evN\ hold and \evU\ fails at a non-valid pair $(\sigma,r)$ with witness $\bar u$. As in the first part, $\sigma$ is a slot, $r\in A(D)$ and $Y_\sigma(D)=Y^*_\sigma$, so $\singleEq(\sigma,r,\bar u)$ is a $D$-feature. If $r\in\singleskel(T^*)$ has rigid root $g$, choose $\theta(M)$ with a symbol root $\neq g$; then $((T^*\theta)|_\sigma)[\bar u]$ has that root (substituting for free indices keeps a symbol root), while $(T^*\theta)|_r$ has root $g$. If $r$ carries an occurrence of $M'\neq M$ (of the same sort), choose $\theta(M)$ and $\theta(M')$ with different symbol roots. If $r$ carries $M(\bar t_r)$, then $\bar t_r\neq\bar t_\sigma[\bar u]$ because the pair is not valid; pick $m$ with $t_{r,m}\neq t_{\sigma,m}[\bar u]$ and put $\theta(M):=\lambda\bar z.z_m$ (term-valued) or $\lambda\bar z.R(z_m,\dots,z_m)$ (formula-valued). The two sides then differ in that argument.\qedhere
\end{itemize}
\end{proof}

\begin{proof}[Proof of \Cref{cor:single:special}]
\emph{(a)}
Every occurrence of a $0$-ary metavariable is a pattern occurrence, and $s_i|_\sigma=\beta^i_M$; so \evRs\ is \evR. \evN\ is vacuous. $Y^*_\sigma=\emptyset$, so \evU\ at $(\sigma,r)$ asks whether $s_i|_r=\beta^i_M$ for all $i$. If $r\in\singleskel(T^*)$, its root is fixed, and \evR\ excludes the equation. If $r$ carries $M'\neq M$ (necessarily of the same sort), the equation is the failure of \evD\ for $M,M'$. If $r$ carries $M$, the pair is valid. So \evRs$\wedge$\evN$\wedge$\evU\ $\Leftrightarrow$ \evR$\wedge$\evD, and \Cref{thm:single:anchor} applies.

\emph{(b)} Formula bodies have symbol roots, and they survive plugging (\Cref{lem:setting:subst}(iii)); so the root of $s_i|_\sigma$ is the root of $\beta^i_M$ at every occurrence: \evRs$\Leftrightarrow$\evR. Assume \evR\ and \evN. For $r\in\singleskel(T^*)$ (of formula sort) a coincidence would make all roots of the $\beta^i_M$ equal to the rigid root at $r$, contradicting \evR. For $r$ an occurrence of the same $M$, $\beta^i_M[\bar t_r]=\beta^i_M[\bar t_\sigma[\bar u]]$ for all $i$ forces $\bar t_r=\bar t_\sigma[\bar u]$ by \evN\ and \Cref{lem:setting:subst}(ii); so either the pair is valid or there is no coincidence. All occurrences have formula sort. What remains of \evU\ is (D$^+$).

\emph{(c)} $T_{\Ind}$ and every schema of \Cref{tab:zf:forms} have a single formula metavariable $P$. In the ZF schemas every occurrence of $P$ is a pattern occurrence (for example $P(x,y,A)$ and $P(x,u,A)$ in $\ReplS$ apply $P$ to distinct bound variables), while $T_{\Ind}$ has the non-pattern occurrences $P(0)$ and $P(Sx)$; all are in $\DT$. Part (b) gives \evR$\wedge$\evN, where \evN\ splits into one condition per argument place of $P$; for $T_{\Ind}$ it says that some motive has $x$ free. A ground template $T^*$ has $C(\{T^*\})=$ all positions and no slot, so $\Feat(\{T^*\})=\{T^*\}$ by \Cref{thm:single:feat}. One instance of a schema never satisfies \evR; two instances whose bodies have different main symbols and together use every argument place do (for $\ReplS$, for example, the bodies $y=x$ and $\neg(x=A)$). If no datum of $T_{\Sep}$ mentions $z$, the learner stays at the strictly more special template with $P(x,z)$ replaced by $P'(x)$, which is sound.
\end{proof}

\paragraph{Weaker events.} The events of \Cref{thm:single:anchor} cannot be weakened to the candidates suggested by first-order patterns.

\begin{proposition}[Weaker events do not suffice]\label{prop:single:weaker}\status{proved; computed}\src{single Prop D.4, D.5, Ex.~D.6}
\leavevmode
\begin{enumerate}[label=(\alph*)]
\item \evR+\evN+\evD: $T^*=\forall x(0=f(x))$, $D=\{\forall x(0=0),\forall x(0=x)\}$. \evRs\ and \evN\ hold and \evD\ is vacuous, but $\forall x(f(0)=f(x))$ covers $D$ and misses $\forall x(0=Sx)$.
\item \evU\ for pattern occurrences only: $T^*=\forall x\forall y(Sf(x,y)=0)\wedge\forall x\forall y(f(x,Sy)=0)$ with the values $\lambda z_1z_2.z_1$, $\lambda z_1z_2.z_2$, $\lambda z_1z_2.Sz_1$. \evRs, \evN, \evD\ hold, and \evU\ holds for every pair whose first component is the pattern occurrence. But the content at the rigid $S$ above $f(x,y)$ equals the content at the derived occurrence $f(x,Sy)$ under $x\mapsto Sx$, so $\forall x\forall y(f'(y,Sx)=0)\wedge\forall x\forall y(f'(y,x)=0)$ covers $D$ and misses $T^*[f:=\lambda z_1z_2.0]$.
\item Projections: $T^*=\forall x\forall y(f(x,y)=0)\wedge\forall x(x=0)$ with values $\lambda z_1z_2.z_1$, $\lambda z_1z_2.z_2$ satisfies \evRs, \evN, but $\forall x\forall y(f(x,y)=0)\wedge\forall x(f(x,x)=0)$ covers $D$. So \evU\ matters even for one metavariable with pattern occurrences only.
\end{enumerate}
\end{proposition}

\noindent\emph{Proof.} (a) $\forall x(f(0)=f(x))$ is in $\DT$ ($f(x)$ is a pattern occurrence) and covers $\forall x(0=0)$ ($f:=\lambda z.0$) and $\forall x(0=x)$ ($f:=\lambda z.z$, $f(0)=0$). An instance with $f(x)=Sx$ has $f(0)=S0\neq0$. The brute-force minimal covering templates are exactly $\forall x(0=f(x))$ and $\forall x(f(0)=f(x))$.

(b) The data are
\[
\forall x\forall y(Sx=0)\wedge\forall x\forall y(x=0),\quad \forall x\forall y(Sy=0)\wedge\forall x\forall y(Sy=0),\quad \forall x\forall y(SSx=0)\wedge\forall x\forall y(Sx=0).
\]
Write $\sigma$ for the left-hand side of the second conjunct (the derived occurrence $f(x,Sy)$) and $r$ for the left-hand side of the first (the rigid $S$ above $f(x,y)$). The data have $(d|_r,d|_\sigma)=(Sx,x),(Sy,Sy),(SSx,Sx)$, and $d|_r=(d|_\sigma)[x\mapsto Sx,\,y\mapsto y]$ in all three. For the covering template, the pattern occurrence $f'(y,x)$ at $\sigma$ reads $f'=\lambda z_1z_2.z_2$, $\lambda z_1z_2.Sz_1$, $\lambda z_1z_2.Sz_2$, and then $f'(y,Sx)$ gives $Sx$, $Sy$, $SSx$ at $r$. For $f:=\lambda z_1z_2.0$, $T^*$ gives $\forall x\forall y(S0=0)\wedge\forall x\forall y(0=0)$, whereas the covering template forces $f'=\lambda z_1z_2.0$ from the second conjunct and then $0$, not $S0$, in the first. \evU\ holds for the pairs with first component $f(x,y)$ (column $(x,y,Sx)$): the right-hand zeros have the wrong root, and the only coincidence, with the derived occurrence $f(x,Sy)$ under $x\mapsto x$, $y\mapsto Sy$, belongs to a valid pair. Brute force (size $\le18$, arguments $\le2$, arity $\le2$) finds exactly two minimal covering templates, $T^*$ (up to argument order) and the displayed one; the referee's code agrees.\src{single e2, e8; referee t2}

(c) With $d_1=\forall x\forall y(x=0)\wedge\forall x(x=0)$ and $d_2=\forall x\forall y(y=0)\wedge\forall x(x=0)$, the derived occurrence $f(x,x)$ reads $x$ from both values. The rigid $x$ of the second conjunct coincides with the slot's content under $x\mapsto x$, $y\mapsto x$; the instance $f:=\lambda z_1z_2.Sz_1$ of $T^*$ is missed.

On random data \evR$\wedge$\evN$\wedge$\evD\ and \evRs$\wedge$\evN$\wedge$\evD\ were wrong in $27$ and $17$ of $248$ cases (always ``says anchor, is not''; referee: $45$, $29$ of $540$), while \evU\ restricted to pattern occurrences was never wrong, so (b) needs special structure.\src{single e4; referee t3}

\begin{example}[\Rich\ cannot be dropped]\label{ex:single:rich}\status{proved}\src{single Ex.~D.7}
In pure set theory without parameters (signature $\in,=$; no constants or function symbols) a term over holes is a hole, so the only term bodies are projections $\lambda\bar z.z_m$. Hence $T^*=\forall x(f(x)\in x)$ has $\inst(T^*)=\{\forall x(x\in x)\}$, and the single datum $\forall x(x\in x)$ is an anchor (every template covering it contains the whole target), although \evRs\ fails. So without \Rich\ the anchor condition is strictly weaker than \evRs$\wedge$\evN$\wedge$\evU; its exact form is open. With parameters \Rich\ holds, since two distinct parameters are closed terms with different roots.
\end{example}

% =====================================================================
\subsection{Rates}\label{app:single:rates}

\begin{proof}[Proof of \Cref{thm:single:rates}]
(a) Let $N\ge1$. By the ``if'' half of \Cref{thm:single:anchor}, which needs no assumption, ``not an anchor'' is contained in $\neg$\evRs$\ \cup\ \neg$\evN$\ \cup\ ($\evRs$\wedge$\evN$\wedge\neg$\evU$)$. For a ground $T^*$ all three events are empty; consistently, every nonempty $D\subseteq\inst(T^*)$ equals $\{T^*\}$ and is an anchor. In general:
\begin{itemize}
\item $\Pr[\text{all roots at }\sigma\text{ equal}]=\sum_hp_h^N\le p_{\max}^{N-1}\sum_hp_h=(1-\rho_\sigma)^{N-1}$, where $p_h$ is the probability of root $h$ at $\sigma$.
\item $\Pr[z_m\text{ unused by all }\theta_i(M)]=(1-\nu_{M,m})^N$.
\item A coincidence at a non-valid pair $(\sigma,r)$ on $D_N$ means that some $\bar u$ has all $\theta_i\in A_{\sigma,r,\bar u}$. Then each of the $\lfloor N/2\rfloor$ disjoint pairs $(\theta_1,\theta_2),(\theta_3,\theta_4),\dots$ has a common $\bar u$; these pairs are independent, each with probability $1-\kappa_{\sigma,r}$.
\end{itemize}
The union bound gives the first inequality. Each term is at most $(1-\lambda)^{\lfloor N/2\rfloor}$, since $N-1\ge\lfloor N/2\rfloor$ for $N\ge1$; there are $c(T^*)$ terms. Lower bounds: under \Rich\ each single failure event is a non-anchor event (\Cref{thm:single:anchor}, ``only if''). All $N$ roots at $\sigma$ equal the most likely head with probability $(1-\rho_\sigma)^N$; $z_m$ is unused with probability $(1-\nu_{M,m})^N$; and a coincidence at $(\sigma,r)$ has probability $\ge\chi_{\sigma,r}^N$ (\Cref{prop:single:kappa}(i)).

(b) $\rho_\sigma>0$ iff two support elements have different roots at $\sigma$; $\nu_{M,m}>0$ iff some support element uses $z_m$; and $\kappa_{\sigma,r}>0$ iff some pair of support elements has no common $\bar u$, which by \Cref{lem:single:basic}(d) (a pairwise condition; finiteness is not used) holds iff the support has no coincidence at $(\sigma,r)$. So $\lambda>0$ iff the support satisfies the three events. Then $c(T^*)(1-\lambda)^{\lfloor N/2\rfloor}\le\bar c\,e^{-\lambda\lfloor N/2\rfloor}\le\bar c\,e^{-\lambda(N-1)/2}$, which is $\le\delta$ iff $N\ge1+(2/\lambda)\ln(\bar c/\delta)$. For a ground $T^*$ the failure probability is $0$ for $N\ge1$.

(c) The version space over tagged calculi factors over tags \citep{claude2026whatfollows}, so the verifier is exact iff the tag-$i$ data are an anchor for every $i$, and
\[
\Pr[\text{not exact}]\le\sum_i\Pr[\text{the tag-}i\text{ data are not an anchor}].
\] Condition on the number $n_i\sim\mathrm{Bin}(N,\pi_i)$ of tag-$i$ data. For all $n_i\ge0$,
\[
\Pr[\text{not an anchor}\mid n_i]\le\bar c_i(1-\lambda_i)^{\lfloor n_i/2\rfloor}:
\]
for $n_i\ge1$ this is (a); for $n_i=0$ the right side is $\bar c_i\ge1$; for a ground tag the left side is $[n_i=0]$ and the right side is $0^{\lfloor n_i/2\rfloor}\ge[n_i=0]$ (with $0^0=1$). Next, $(1-\lambda)^{\lfloor n/2\rfloor}\le(1-\lambda)^{(n-1)/2}\le e^{\lambda/2}e^{-\lambda n/2}\le\sqrt e\,e^{-\lambda n/2}$, as $\lambda\le1$ (for $\lambda=1$ check $n=0,1$ directly). And $\mathbb E[e^{-an_i}]=(1-\pi_i(1-e^{-a}))^N\le\exp(-N\pi_i(1-e^{-a}))$, with $1-e^{-a}\ge3a/4$ for $a=\lambda_i/2\le1/2$ (the left side is concave, and the inequality holds at $a=1/2$). This gives $\sqrt e\,\bar c_i\exp(-3N\pi_i\lambda_i/8)$ per tag. Lower bounds: with probability $(1-\pi_i)^N$ there is no tag-$i$ datum, and then $\Acc$ of the tag-$i$ data is empty; this is the exact failure probability of a ground tag. Under \Rich, all tag-$i$ data share the head $h$ at the pattern occurrence, or there are none, with probability $(1-\pi_i+\pi_i(1-\rho))^N=(1-\pi_i\rho)^N$, and then they are no anchor.
\end{proof}

The \evU-terms vanish for a single formula metavariable because, under \evR$\wedge$\evN, \evU\ holds automatically (\Cref{cor:single:special}(b)), and all occurrence events coincide with \evR; so the bound is $(1-\rho)^{N-1}+\sum_m(1-\nu_m)^N$.

\begin{proposition}[Coincidence terms]\label{prop:single:kappa}\status{proved; computed}\src{single Prop E.2}
For $(\sigma,r)\in U(T^*)$: (i) the probability of a coincidence at $(\sigma,r)$ on $D_N$ is $\ge\chi_{\sigma,r}^N$, and under \Rich\ such a coincidence makes $D_N$ a non-anchor; (ii) $\chi_{\sigma,r}^2\le1-\kappa_{\sigma,r}$; (iii) if $\Lambda$ has finite support and $\mathcal K$ is the family of maximal support sets lying in a common $A_{\sigma,r,\bar u}$, the coincidence probability equals $\sum_{\emptyset\neq\mathcal J\subseteq\mathcal K}(-1)^{|\mathcal J|+1}\Lambda(\bigcap\mathcal J)^N\in[\chi^N,|\mathcal K|\chi^N]$. So its exact rate is $\chi$, which can be smaller than the rate $\sqrt{1-\kappa}$ of the upper bound; they agree if $|\mathcal K|=1$.
\end{proposition}

\begin{proof}[Proof of \Cref{prop:single:kappa}]
(i) For each $\bar u$, $\Pr[\text{all }\theta_i\in A_{\bar u}]=\Lambda(A_{\bar u})^N$; take the supremum. Under \Rich, a coincidence at a non-valid pair is a failure of \evU, so $D_N$ is not an anchor (by \Cref{thm:single:anchor} ``only if'' if \evRs\ and \evN\ hold, and also otherwise). (ii) For every $\bar u$, $1-\kappa=\Lambda\otimes\Lambda\{(\theta,\theta'):\exists\bar u'\ \theta,\theta'\in A_{\bar u'}\}\ge\Lambda(A_{\bar u})^2$. (iii) The event says that the set of sampled support points is compatible, i.e.\ lies in some $A_{\bar u}$, i.e.\ in some $K\in\mathcal K$; inclusion--exclusion over $\mathcal K$ gives the formula. Each $K$ lies in some $A_{\bar u}$, so $\Lambda(K)\le\chi$ and the probability is at most $|\mathcal K|\chi^N$. Conversely $A_{\bar u}\cap\operatorname{supp}\Lambda$ is compatible for every $\bar u$, hence contained in a maximal compatible set; so $\chi=\max_K\Lambda(K)$, and the probability is at least $\chi^N$. If $\mathcal K=\{K\}$ is a singleton, a pair (or a single point) is compatible iff it lies in $K$, so $1-\kappa=\Lambda(K)^2=\chi^2$.
\end{proof}

\paragraph{Two worked laws.} \emph{First law} (the referee's example): $T^*=\forall x(0=f(x))$, with $f$ equal to $\lambda z.0$, $\lambda z.z$, $\lambda z.Sz$ with probabilities $0.4$, $0.3$, $0.3$. The only non-valid pair is (slot of $f(x)$, rigid $0$), and its compatible set is $\{\lambda z.0,\lambda z.z\}$ (with $\bar u:x\mapsto0$); $\lambda z.Sz$ is in no $A_{\bar u}$. \evRs\ fails iff all values are equal, \evN\ fails iff all values are $\lambda z.0$, and \evU\ fails iff all values lie in $\{0,z\}$. So $\Pr[\text{not an anchor}]=0.7^N+0.3^N$, which at $N=8$ is $0.057714$. It lies between the coincidence lower bound $\chi^8=0.7^8=0.057648$ and the upper bound $0.4^7+0.4^8+0.49^4=0.059942$ of \Cref{thm:single:rates}(a); the root and non-vacuity lower bounds alone give $0.4^8=0.000655$. \emph{Second law:} $T^*=\forall x\forall y(f(x)=g(y))$ with $f\sim U\{z,Sz\}$ and $g\sim U\{z,0,Sz\}$ independent; write $\theta$ as the pair (value of $f$, value of $g$). At the pair $(f(x),g(y))$ the maximal compatible sets are $\{(z,z),(Sz,Sz)\}$ (with $\bar u:x\mapsto y$), $\{(z,0)\}$ ($x\mapsto0$) and $\{(z,Sz)\}$ ($x\mapsto Sy$), of masses $1/3,1/6,1/6$; so $|\mathcal K|=3$, $\chi=1/3$, and $1-\kappa=1/9+2/36=1/6$, so $\sqrt{1-\kappa}=0.408>\chi$. At $N=8$ the exact coincidence probability is $(1/3)^8+2(1/6)^8=1.54\cdot10^{-4}$, against $\chi^N=1.52\cdot10^{-4}$ and $(1-\kappa)^4=7.7\cdot10^{-4}$. \emph{Non-circular check of \Cref{thm:single:rates}(a):} the referee computed exact non-anchor probabilities by inclusion--exclusion over the support, deciding anchor status by feature search rather than by the events of \Cref{thm:single:anchor}, for induction, Separation and the template of \Cref{prop:single:weaker}(b); the bound held at every $N\in\{2,4,8,12,16,24\}$ (for induction at $N=24$: exact $0.00100$, bound $0.00535$).\src{referee t7}  The computations of Proposition~\ref{prop:single:kappa}(iii) were checked on every non-valid pair of five laws, and exact non-anchor probabilities for two small laws were cross-checked against an independent feature-template test with no disagreement.\src{single e10}

% =====================================================================
\subsection{Escalations}\label{app:single:esc}

Let $\Pi(D):=(C(D),(Y_\sigma(D))_\sigma,E(D))$, where $E(D)$ is the set of pairs $(\sigma,r)$ carrying an Eq-feature. As data grow, $C$ shrinks, scopes grow, and while both are fixed $E$ shrinks; and by \Cref{thm:single:feat} a query outside $\Acc(D)$ changes $\Pi$ (\Cref{lem:single:monotone}).

\begin{lemma}[Monotonicity]\label{lem:single:monotone}\status{proved}\src{single Lemma F.1}
Let $D\subseteq D'$ be nonempty.
(i) $C(D')\subseteq C(D)$ and $A(D')\subseteq A(D)$; a slot of $D$ that is still in $A(D')$ is a slot of $D'$; and $Y_\sigma(D)\subseteq Y_\sigma(D')$ whenever $\sigma$ is a slot of both.
(ii) If $C(D')=C(D)$ and $Y(D')=Y(D)$, then $E(D')\subseteq E(D)$, with the same maps $\bar u$.
(iii) If $q\notin\Acc(D)$, then $\Pi(D\cup\{q\})\neq\Pi(D)$.
\end{lemma}

\begin{proof}
(i) is immediate from the definitions: a slot of $D'$ has its parent in $C(D')\subseteq C(D)$, and scopes are unions over the data. (ii) A feature of $D'$ restricts to $D$, and its $\bar u$ is forced on $Y_\sigma(D)=Y_\sigma(D')$ (\Cref{lem:single:basic}(d)). (iii) If $\Pi$ is unchanged, then $q$ has every Sym feature ($C$ unchanged), every Scope feature ($Y$ unchanged) and every Eq feature ($E$ unchanged, with $\bar u$ forced by $D$). So $q\in\Feat(D)=\Acc(D)$ by \Cref{thm:single:feat}.
\end{proof}

\begin{proof}[Proof of \Cref{thm:single:esc}]
Let $D_j=\{d,q_2,\dots,q_j\}$. By \Cref{lem:single:monotone}, each step $D_{j-1}\to D_j$ changes $\Pi$, and then it does one of the following: ($\alpha$) shrink $C$; ($\beta$) keep $C$ and enlarge some scope (scopes only grow, and slots are unchanged when $C$ is); ($\gamma$) keep $C$ and every scope, and remove at least one pair from $E$ (by (ii), $E$ can only shrink then). ($\alpha$) happens at most $|C(\{d\})|=n$ times. A position $\sigma$ is a slot during one interval of steps: it becomes a slot when it leaves $C$, and it leaves $A$ for ever when its parent leaves $C$. During that interval $Y_\sigma$ only grows, inside $\{0,\dots,\singlebd(\sigma)-1\}$; so ($\beta$) happens at most $\sum_\sigma\singlebd(\sigma)$ times. A pair $(\sigma,r)$ can lie in $E$ only while $\sigma$ is a slot and $r\in A$, which is a single interval. Once removed, it never returns: if it were in $E$ at times $t_1<t_3$ but not at $t_2\in(t_1,t_3)$, then $\sigma$ is a slot and $r\in A$ at $t_2$, and the feature at $t_3$ restricts to the data at $t_2$ (\Cref{lem:single:basic}(d)), a contradiction. So ($\gamma$) happens at most once per ordered pair of incomparable positions of $d$. The second inequality uses $\singlebd(p)\le\singlebd_{\max}$ and at most $n(n-1)$ ordered pairs. From no data the first query is escalated, and $\singlebd(p)\le n-1$, so $\singleEsc(\DT;N)\le1+N+N(N-1)+N(N-1)\le2N^2+1$.

\emph{The linear lower bound.} First-order targets already give $\singleEsc(\DT;N)\ge N+1$ for $N\ge5$. With $k=N-3$, the chain $S^k0=0$, $S^ka=0$, $S^{k-1}0=0$, \dots, $0=0$, $0=S0$, $\neg(0=0)$ has $N+1$ members of size $\le N$, and the successive acceptance sets are $\{S^k0=0\}$, $\{S^kt=0\}$, $\{S^{k-1}t=0\}$, \dots, $\{t=0\}$, $\{t=t'\}$ ($t,t'$ closed): each query lies outside the previous set. All are instances of the universal template. Checked for $N=5,8,12,16$.\src{single e8}
\end{proof}

The quadratic term is real:

\begin{proposition}[Quadratic escalations]\label{prop:single:quadratic}\status{proved; computed}\src{single Prop F.8}
Let $d_1=\forall y_1\dots\forall y_m(0=0\wedge\dots\wedge0=0)$ with $m$ atoms, $d_2$ the sentence $d_1$ with every left-hand side replaced by a parameter $a$, and $q_{j,k}$ ($j\le m$, $k<m$) the sentence $d_1$ with the $j$-th left-hand side replaced by the bound variable of index $k$.
(i) In the order $d_1,d_2,q_{1,0},\dots,q_{1,m-1},q_{2,0},\dots,q_{m,m-1}$, each sentence lies outside the acceptance set of its predecessors ($d_2$ violates a Sym feature, each $q_{j,k}$ a Scope feature). All have size $5m-1$. Hence $2+\lfloor(N+1)/5\rfloor^2\le\singleEsc(\DT;N)\le2N^2+1$ for $N\ge4$: $\singleEsc(\DT;N)=\Theta(N^2)$.
(ii) $\Ind(d_1),\Ind(d_2),\Ind(q_{j,k})$, in the same order, are instances of $T_{\Ind}$ of size $20m+1$, and each lies outside the acceptance set of its predecessors.
\end{proposition}

\noindent The earlier analysis conjectured that escalations are at most \emph{linear} in the size of the first escalated datum\src{prior induction C11}. By (ii), $1+m^2$ escalations follow a first datum of size $n=20m+1$, which exceeds $c\,n$ once $m>20c+1$. So this conjecture is false: for $\DT$, for the class with nested arguments (whose acceptance sets are smaller), for the universal template and for the induction target. The refutation is asymptotic (for $m=2,\dots,5$ the counts $5,10,17,26$ are below $n=41,\dots,101$). The quadratic bound holds and is tight. Its quadratic part comes from scope growth, which needs deep binder nesting (\Cref{prop:single:binderfree}).

\begin{proof}[Proof of \Cref{prop:single:quadratic}]
(i) $d_1$ and $d_2$ differ exactly at the $m$ left-hand sides (roots $0$ and $a$). So $C(\{d_1,d_2\})$ is everything else, and the left-hand sides are slots with empty scope. Each $q_{j,k}$ agrees with $d_1$ off the slots, so $C$ does not change along the chain. At slot $j$, $q_{j,k}$ has the bound variable with index $k$ free (relative to the slot, which lies under the $m$ binders and no others). Earlier data have at slot $j$ only $0$, $a$, or other bound variables $k'<k$. So $q_{j,k}$ violates $\singleScope(\text{slot }j)$ of its predecessors, and $q_{j,k}\notin\Feat=\Acc$. (Without the hard half of \Cref{thm:single:feat}: $T_0$ of the predecessors is a covering $\DT$ template that misses $q_{j,k}$.) The size is $m$ binders, $m-1$ conjunctions and $3m$ symbols in the atoms: $5m-1$. The chain has $2+m^2$ members. For $N\ge4$, $m=\lfloor(N+1)/5\rfloor\ge1$ gives size $5m-1\le N$; the upper bound is \Cref{thm:single:esc}.

(ii) The motive occurs four times in $\Ind(\varphi)$: as $\varphi(0)$, $\varphi(x)$, $\varphi(Sx)$ and in the conclusion; since $x$ does not occur, all four copies are $\varphi$. After $\Ind(d_1)$ and $\Ind(d_2)$ the left-hand sides in all four copies are slots, and their scopes contain no index bound by the motive's own $m$ binders: $x$ does not occur, and $d_1,d_2$ have no indices at the slots. $\Ind(q_{j,k})$ puts the index $k$ (bound by the motive's binders) at slot $j$ of each copy, which no predecessor has. The size is $4(5m-1)+5=20m+1$. For $m>20c+1$, $1+m^2>c(20m+1)$.
\end{proof}

\emph{Computed.} Chains of length $6,18,38,66$ for $m=2,4,6,8$; for the universal template and $N\in\{4,5,9,14,19,23,24,29,34\}$ exactly $2+\lfloor(N+1)/5\rfloor^2$ escalations with sizes $\le N$; for the induction target, first data of size $41,61,81,101$ followed by $5,10,17,26$ escalations. Every step was certified twice, by the feature verifier and by an explicit covering $T_0$ that misses the query; the referee's independent code agrees.\src{single e8, e12; referee t5, t5b}

\begin{proposition}[Binder-free data]\label{prop:single:binderfree}\status{proved; computed}\src{single Prop F.9}
For a chain as in \Cref{thm:single:esc}, at most $n-1$ steps of kind ($\gamma$) remove a pair whose source has empty scope. If $d$ has no binder, $m-1\le2n-1$; on binder-free sentences $N+1\le\singleEsc\le2N$ ($N\ge5$).
\end{proposition}

\begin{proof}[Proof of \Cref{prop:single:binderfree}]
(i) For $D\ni d$ define $p\sim_Dp'$ iff $p=p'$, or $p,p'$ are positions of every datum and $d'|_p=d'|_{p'}$ for all $d'\in D$. This is an equivalence relation on the positions of $d$, and it refines as $D$ grows. Consider a ($\gamma$)-step $D\to D'=D\cup\{q\}$; it changes neither $C$ nor the scopes, hence neither the slots nor $A$. Suppose it removes $(\sigma,r)$ with $Y_\sigma(D)=\emptyset$. Then $\singleEq(\sigma,r,\emptyset)$ was a $D$-feature, i.e.\ $d'|_r=d'|_\sigma$ for all $d'\in D$, so $\sigma\sim_Dr$. It is not a $D'$-feature, although $\sigma$ is still a slot, $r$ still admissible and $\FV(q|_\sigma)\subseteq Y_\sigma(D')=\emptyset$; so $q|_r\neq q|_\sigma$ and $\sigma\not\sim_{D'}r$. This is a strict refinement. An equivalence relation on $n$ points refines strictly at most $n-1$ times.

(ii) If $d$ has no binder, every slot $\sigma$ has $\singlebd(\sigma)=0$, since its ancestors lie in $C$ and carry the symbols of $d$. So all scopes are empty: there are no ($\beta$)-steps, and every ($\gamma$)-step is of the kind in (i). With at most $n$ ($\alpha$)-steps, $m-1\le n+(n-1)$. From no data add $1$ for the first escalation; the lower bound $N+1$ is the binder-free chain above.
\end{proof}

\emph{Computed.} $40$ greedy adversarial chains from binder-free first data of $8$--$18$ positions: no chain exceeded $2n-1$, and every ($\gamma$)-step strictly refined $\sim$.\src{single e13}

\begin{lemma}[Triangle constraints on equation kills]\label{lem:single:triangle}\status{proved; computed}\src{single Lemma F.10}
Let $D\subseteq D'$ with $C(D')=C(D)$ and $Y(D')=Y(D)$, and write $\bar u_{ab}$ for the forced map of a pair $(a,b)\in E$.
(a) If $(k,i),(i,j)\in E(D')$, $k\perp j$, and $k,j$ have the same sort, then $(k,j)\in E(D')$ with $\bar u_{kj}=\bar u_{ki}[\bar u_{ij}]$ on $Y_k$.
(b) If $(k,i),(k,j),(i,j)\in E(D)$ and $(k,i),(k,j)\in E(D')$, then $(i,j)\in E(D')$.
(c) If $(i,k),(k,j),(i,j)\in E(D)$ and $(i,k),(k,j)\in E(D')$, then $(i,j)\in E(D')$.
Hence a ($\gamma$)-step that removes $(i,j)$ also removes one of $(k,i),(k,j)$ for every common in-neighbour $k$, and one of $(i,k),(k,j)$ for every intermediate $k$.
\end{lemma}

\begin{proof}
First, $\bar u_{ki}[\bar u_{ij}]$ is defined on $Y_k$: for $y\in Y_k$, $\bar u_{ki}(y)$ is read off a datum $d$ with $y\in\FV(d|_k)$, so its free indices lie in $\FV(d|_i)\subseteq Y_i$, the domain of $\bar u_{ij}$. (a) $k$ is a slot and $j\in A$. For $d'\in D'$, $d'|_j=(d'|_i)[\bar u_{ij}]=((d'|_k)[\bar u_{ki}])[\bar u_{ij}]=(d'|_k)[\bar u_{ki}[\bar u_{ij}]]$, and $\FV(d'|_k)\subseteq Y_k$. (b) By (a) on $D$ and forcing, $\bar u_{kj}=\bar u_{ki}[\bar u_{ij}]$ on $Y_k$. Let $q\in D'\setminus D$. The maps of $(k,i)$ and $(k,j)$ on $D'$ are the forced maps on $D$ (the domains are unchanged). So $q|_i=(q|_k)[\bar u_{ki}]$ and $q|_j=(q|_k)[\bar u_{kj}]=((q|_k)[\bar u_{ki}])[\bar u_{ij}]=(q|_i)[\bar u_{ij}]$; moreover $\FV(q|_i)\subseteq Y_i(D')=Y_i(D)$. So $\singleEq(i,j,\bar u_{ij})$ holds on $D'$. (c) Apply (a) to $D'$; its map $\bar u_{ik}[\bar u_{kj}]$ equals $\bar u_{ij}$ on $D$ by forcing.
\end{proof}

\begin{conjecture}[Equation kills are linear]\label{conj:single:kill}\status{conjecture}\src{single Conj.~F.7}
The number of ($\gamma$)-steps is $O(n)$; hence $\singleEsc(\DT;N)=O(N(1+\singlebd_{\max}))$ for binder depth $\le\singlebd_{\max}$.
\end{conjecture}

\noindent The evidence is weak: greedy adversarial searches found at most about $1.4m$ kill steps on $m\le12$ slots, but they also missed the scope-growth chains of \Cref{prop:single:quadratic}.\src{single e6, e6b} Surviving equations need not form an equivalence relation; they compose like a preorder and obey a triangle rule (\Cref{lem:single:triangle}), and chains of preorders on $m$ points can have length of order $m^2$. \Cref{prop:single:binderfree} proves the conjecture when all scopes are empty.

\emph{Where the difficulty of \Cref{conj:single:kill} sits.} In a phase with $C$ and the scopes fixed, ``same value in every datum'' is an equivalence relation on the nodes $A(D)\cup\{(\sigma,\bar u_{\sigma r})\}$ (the value of $(\sigma,\bar u)$ in $d$ being $(d|_\sigma)[\bar u]$); it refines with each datum, and $(\sigma,r)$ survives iff $(\sigma,\bar u_{\sigma r})\equiv r$. So a phase has fewer ($\gamma$)-steps than nodes: linearly many if every slot uses $O(1)$ distinct maps, but there can be order $n^2$ nodes, and up to $n+\sum\singlebd$ phases. \emph{Computed:} $587$ composition instances and $935$ instances each of (b) and (c), $240$ of them in steps that removed some pair, with no failure.\src{single e13}

% =====================================================================
\subsection{Unions}\label{app:single:unions}

\begin{proof}[Proof of \Cref{cor:single:unions}]
(i) If $w_0,w_1,\dots$ and $L_n=\inst(T_n)$ witness infinite elasticity, then $w_n\notin\Acc(\{w_0,\dots,w_{n-1}\})$ for $n\ge1$, so the sequence is bounded by \Cref{thm:single:esc} with $d=w_0$. For thickness: $P:=\lambda z.(0=0)$ puts $\forall x(0=0)\wedge(0=0)$ in $\inst(\forall xP(x)\wedge P(t))$ for every closed $t$, and $P:=\lambda z.(z=z)$ gives $\forall x(x=x)\wedge(t=t)$, which lies in the $t$-template only (matching forces $P=\lambda z.(z=z)$ from the pattern occurrence). (ii) is the cited preservation result together with ``finite elasticity gives anchors'' and the eventual exactness of the cautious verifier once an anchor is in the data \citep{claude2026whatfollows}.

(iii) is \Cref{prop:single:union-esc}.
\end{proof}

\begin{proposition}[Escalations for unions]\label{prop:single:union-esc}\status{proved; computed}\src{single Cor F.4}
With a $k$-ary relation symbol, a unary function symbol and a constant, $\singleEsc(H_k(\DT);N)\ge\lfloor(N-1)/k\rfloor^k$; over $=,+,S,0$, $\singleEsc(H_k(\DT);N)\ge(\lfloor(N-1)/k\rfloor-1)^k$ for $N\ge2k+1$, already for one target.
\end{proposition}

\begin{proof}
With a $k$-ary relation symbol $P$, a unary function symbol $g$ and a constant $c$, the chain of \citet{claude2026whatfollows} for the term signature transfers verbatim: query $P(g^{a_1}c,\dots,g^{a_k}c)$ for $a\in\{0,\dots,n\}^k$, $n+1=\lfloor(N-1)/k\rfloor$, by non-increasing $\sum_ja_j$. Over $=,+,S,0$, let $n+1:=\lfloor(N-1)/k\rfloor-1\ge1$ (so $N\ge2k+1$), and query
\[
q_a:=\bigl(S^{a_1}0+(S^{a_2}0+\dots+S^{a_k}0)\bigr)=0,\qquad a\in\{0,\dots,n\}^k,
\]
ordered by non-increasing $\sum_ja_j$. Each has size $k+1+\sum_j(a_j+1)\le k+1+k(n+1)\le N$, and all are instances of the target $(x_1+\dots+x_k)=0$ with $0$-ary term metavariables $x_j$. Every earlier query $a'$ has $a'_j\ge a_j+1$ for some $j$; otherwise $a'\le a$ componentwise with $a'\neq a$, so $\sum a'<\sum a$, contradicting the order. Hence the earlier queries lie in $\bigcup_j\inst(T_j)$ with $T_j:=(x_1+\dots+S^{a_j+1}(x_j)+\dots+x_k)=0\in\DT$, and this union misses $q_a$. So $q_a\notin\Acc_k(\text{earlier queries})$; the first query is escalated because $\Acc_k(\emptyset)=\emptyset$. This gives $(n+1)^k$ escalations. With a $k$-ary $P$ the same proof works with overhead $1$ instead of $k+1$.
\end{proof}

\emph{Computed.} For $(k,n+1)=(1,6),(2,3),(2,4),(3,2)$: $6,9,16,8$ escalations on sentences of size $\le N=8,9,11,10$, each step certified by the explicit $k$-union and, for the first $12$ steps, by the brute-force union verifier.\src{single e16} The encoded bound over $=,+,S,0$ was made explicit after the referee round, so it was not checked by the independent referee (\Cref{app:ver:unref}).

\paragraph{Explicit anchors for unions.} Let the target be $\bigcup_{i\le k'}\inst(T^*_i)$, $k'\le k$, $T^*_i\in\DT$, and $\Theta_i=\{\theta:T^*_i\theta\in D\}$. Let $\singleFail(T^*)$ consist of the failure sets of the three events: $\{\theta:\singleroot((T^*\theta)|_\sigma)=h\}$, $\{\theta:z_m\notin\theta(M)\}$, and $\{\theta:(T^*\theta)|_r=((T^*\theta)|_\sigma)[\bar u]\}$ for non-valid $(\sigma,r)$.

\begin{proposition}[Explicit untagged anchors]\label{prop:single:untagged-anchor}\status{proved; computed}\src{single Prop F.5}
If for each $i$ the set $\Theta_i$ is nonempty and not covered by $k$ members of $\singleFail(T^*_i)$, then $D$ is an anchor for the union in $H_k(\DT)$. For a ground $T^*_i$ (a single axiom) $\singleFail(T^*_i)=\emptyset$, and the condition says that the axiom itself is a datum.
\end{proposition}

\noindent Nonemptiness matters for ZFC's single axioms: with target $\inst(T_{\Ind})\cup\{\Ext\}$ and $k=2$, the union verifier rejects Extensionality until it is a datum (computed below; the repair that added this hypothesis is recorded in \Cref{app:ver:corr}). For one formula metavariable, $\singleFail$ reduces to ``all roots equal $f$'' and ``argument $m$ unused'', the failure family of the earlier untagged PA analysis (\Cref{cor:many:failure}).

\begin{proof}[Proof of \Cref{prop:single:untagged-anchor}]
Let $\bigcup_{l\le k}\inst(T_l)\supseteq D$, $T_l\in\DT$. Assign each $\theta\in\Theta_i$ to some $l$ with $T^*_i\theta\in\inst(T_l)$ (matchers are unique, so data and substitutions correspond). This splits $\Theta_i$ into at most $k$ parts, and into at least one nonempty part, since $\Theta_i\neq\emptyset$. Suppose no part's data were an anchor for $\inst(T^*_i)$. By the ``if'' half of \Cref{thm:single:anchor} (which needs no \Rich), each nonempty part then fails \evRs, \evN\ or \evU, i.e.\ lies in a member of $\singleFail(T^*_i)$ (all roots at $\sigma$ equal some $h$; or $z_m$ unused; or, with \evRs\ and \evN\ holding, a coincidence $\bar u$ at a non-valid pair). For a ground $T^*_i$ a nonempty part is $\{T^*_i\}=\inst(T^*_i)$, an anchor. So $\Theta_i$ would be covered by at most $k$ members of $\singleFail(T^*_i)$, contrary to the hypothesis. Hence some part, contained in $\inst(T_l)$, is an anchor, and $\inst(T_l)\supseteq\inst(T^*_i)$. This holds for every $i$.
\end{proof}

\noindent\emph{Computed.} With target $\inst(T_{\Ind})\cup\{\Ext\}$ and $k=2$, three induction data and no Extensionality datum, the union verifier rejects Extensionality; with Extensionality among the data it accepts it; fresh induction instances are accepted in both cases. So the hypothesis $\Theta_i\neq\emptyset$ is needed.\src{single e15}

\begin{proposition}[The union verifier]\label{prop:single:union-verifier}\status{proved; computed}\src{single Prop F.6}
$q\in\Acc_k(D)$ iff every partition of $D$ into at most $k$ nonempty blocks has a block $B$ with $q\in\Acc(B)$. So $\Acc_k$ is decidable in time (number of such partitions)$\times$poly.
\end{proposition}

\noindent This is the case $N=\emptyset$ of \Cref{lem:many:cautious}(b), in the form \eqref{eq:many:partition}; \Cref{thm:many:complexity} shows that the test is polynomial for $k\le2$ but coNP-complete for each fixed $k\ge3$.

\begin{proof}[Proof of \Cref{prop:single:union-verifier}]
A $k$-union $\bigcup_{l\le k}\inst(T_l)$ contains $D$ iff some partition of $D$ into at most $k$ nonempty blocks $B_1,\dots,B_j$ has each block covered by its own $T_l$ ($l\le j$), the remaining templates being unconstrained. Suppose that for some such partition $q\notin\Acc(B_l)$ for every $l\le j$. Choose $T_l$ covering $B_l$ with $q\notin\inst(T_l)$, and for $l>j$ a ground template other than $q$ (possible since $\Acc(\emptyset)=\emptyset$). This union contains $D$ and misses $q$, so $q\notin\Acc_k(D)$. Conversely, if every partition has a block $B$ with $q\in\Acc(B)$, let $U=\bigcup_l\inst(T_l)\supseteq D$; assigning each datum to a template that covers it gives a partition, some block $B$ of it has $q\in\Acc(B)\subseteq\inst(T_l)\subseteq U$. Each test $q\in\Acc(B)$ is polynomial by \Cref{cor:single:poly}.
\end{proof}

\noindent\emph{Computed.} On unlabelled data from PA induction and $\in$-induction (three instances each), $k=1$ accepts the false $(0=0\wedge\forall x(x=x\to Sx=Sx))\to\forall x\,x=0$, while $k=2,3$ accept all five fresh instances tried and reject it and two other non-instances (reproduced by the referee).\src{single e9}

% =====================================================================
\subsection{Complexity of the version space}\label{app:single:complexity}

\begin{proof}[Proof of \Cref{prop:single:blowup}]
Each datum has $4+(4n-1)+1=4n+4$ positions. $C(D_n)=\forall x(\square=\square)\wedge C_n$; the two slots have column $(x,0)$ and scope $\{x\}$. The $D_n$-features are the Sym and Scope features, the two slot-to-slot equations (identity map), and, for each of the $2n$ zeros of $C_n$ and each slot, ``content $=$ slot content with $x:=0$''. There are no formula slots, so every metavariable is term-valued with pattern occurrences at the slots, and by \Cref{thm:single:min}(b) a saturated template has at each zero either the rigid $0$ or a derived occurrence. Every covering template is $\gen$ one of the templates $T_A$ ($A$ a set of zeros): $f(x)$ at both slots (merging two metavariables at the slots is a specialization that still covers) and $f(0)$ exactly at the zeros in $A$. For $A\neq B$, the instance $f:=\lambda z.Sz$ of $T_A$ has $S0$ exactly at the positions of $A$, so $T_A$ and $T_B$ have incomparable instance sets. Hence every $T_A$ is minimal ($T\preceq T_A$ covering implies $T\gen T_B\preceq T_A$ for some $B$, forcing $B=A$ and $T\equiv T_A$), and $\Min(D_n)=\{T_A\}$, with $4^n$ elements. By \Cref{thm:single:feat}, $\Acc(D_n)=\{\forall x(t=t)\wedge C_n:\ \FV(t)\subseteq\{x\},\ t[0/x]=0\}=D_n$, since a term with root $S$ or $+$ keeps its root under $[0/x]$, so $t\in\{x,0\}$.
\end{proof}

\noindent\emph{Computed.} $|\Min(D_n)|=4,16,64$ for $n=1,2,3$ from the saturated templates, and $4,16$ for $n=1,2$ by brute force (also by the referee).\src{single e2; referee t2}

\begin{proposition}[Existence of an lgg]\label{prop:single:lgg}\status{proved; computed}\src{single Prop G.3}
$D$ has a least covering template iff (i) no Eq-feature of $D$ has $r\in C(D)$, and (ii) in the graph on $\singleslots(D)$ with an edge $\sigma\to r$ for each slot-to-slot Eq-feature, every weakly connected component $K$ has a \emph{source} $s$ with $s\to x$ for all $x\in K\setminus\{s\}$. Then the lgg is the template $T_{\mathrm{all}}$: $C(D)$ with $M_K(Y_s)$ at each source $s$ and $M_K(\bar u_{s,x})$ at the other slots $x$ of $K$, where $\bar u_{s,x}$ is the forced map. Both conditions, and $T_{\mathrm{all}}$, are computed in polynomial time from the $D$-features (\Cref{cor:single:poly}).
\end{proposition}

\begin{proof}[Proof of \Cref{prop:single:lgg}]
($\Leftarrow$) $T_{\mathrm{all}}\in\DT$, and it covers $D$ with $\theta_d(M_K):=(d|_s)[\bar z/Y_s]$, since $d|_x=(d|_s)[\bar u_{s,x}]$. Let $q=T_{\mathrm{all}}\theta$. It has every Sym feature, and every Scope feature: $\FV(q|_s)\subseteq Y_s$, and $\FV(q|_x)\subseteq\FV(\bar u_{s,x})\subseteq Y_x$, because $\bar u_{s,x}$ is read off data at $x$. Let $\singleEq(\sigma,r,\bar u)$ be a $D$-feature. By (i), $r$ is a slot, and $\sigma,r$ lie in one component $K$, with source $s$ (put $\bar u_{s,s}:=\mathrm{id}$). Then $q|_r=\theta(M_K)[\bar u_{s,r}]$ and $(q|_\sigma)[\bar u]=\theta(M_K)[\bar u_{s,\sigma}[\bar u]]$. On $D$ both $\bar u_{s,r}$ and $\bar u_{s,\sigma}[\bar u]$ solve the equation of the pair $(s,r)$, so they coincide on $Y_s$ by forcing, and $q|_r=(q|_\sigma)[\bar u]$. So $\inst(T_{\mathrm{all}})\subseteq\Feat(D)=\Acc(D)\subseteq\inst(T_{\mathrm{all}})$. \emph{Syntactic leastness:} let $U$ cover $D$ and take a saturated $U'$ with $U\gen U'$ (\Cref{thm:single:min}(a)). An occurrence of $U'$ in $C(D)$ would be a derived occurrence with an Eq-feature landing in $C(D)$, which (i) excludes; so all occurrences of $U'$ are at slots, every slot is one, and the skeleton of $U'$ is $C(D)$. Map each metavariable $L$ of $U'$, with pattern slot $\sigma_L$ and pattern arguments $\bar y$, to $\lambda\bar z.(T_{\mathrm{all}}|_{\sigma_L})[\bar z/\bar y]$. Each other occurrence $L(\bar t)$ at $r$ gives the $D$-feature $(\sigma_L,r,\bar y\mapsto\bar t)$, so $r$ lies in the component of $\sigma_L$, and $T_{\mathrm{all}}|_r=(T_{\mathrm{all}}|_{\sigma_L})[\bar y\mapsto\bar t]$ by forcing. So this substitution maps $U'$ to $T_{\mathrm{all}}$, and $U\gen U'\gen T_{\mathrm{all}}$.

($\Rightarrow$) Let $T_\ell$ be least. $T_\ell\preceq T_0(D)$ makes every position of $C(D)$ rigid in $T_\ell$. Suppose an Eq-feature $\varphi=(\sigma,r,\bar u)$ had $r\in C(D)$. Since $T_\ell$ covers $D$, $\singleskel(T_\ell)\subseteq C(D)$, so the slot $\sigma$ (whose parent is rigid) carries an occurrence $N(\bar t)$ in $T_\ell$. From $T_\ell=T_\varphi\tau$ we get $T_\ell|_\sigma=\tau(M_\sigma)[Y_\sigma]$, so $\tau(M_\sigma)=\lambda\bar z.N(\bar t')$, and then $T_\ell|_r=\tau(M_\sigma)[\bar u]=N(\bar t'[\bar u])$ is an occurrence at a rigid position, a contradiction. This proves (i). So the occurrences of $T_\ell$ are exactly the slots. The same argument shows that a feature $\sigma\to r$ makes $\sigma$ and $r$ occurrences of the same metavariable. For a metavariable $N$ with pattern occurrence at $s$, the pair $(s,x)$ is a feature for every other occurrence $x$ of $N$, because $T_\ell$ covers $D$. So the components are the occurrence sets of the metavariables, each with the source $s$: this is (ii).
\end{proof}

\begin{proof}[Proof of \Cref{cor:single:anchor-lgg}]
By \Cref{thm:single:anchor} (here \Rich\ is used), \evRs, \evN\ and \evU\ hold. So $C(D)=\singleskel(T^*)$, $\singleslots(D)=\Occ(T^*)$, $Y_\sigma(D)=Y^*_\sigma$, and every Eq-feature comes from a valid pair: $r$ is an occurrence of the same metavariable. Thus no feature lands in $C(D)$, which is (i). The slot graph has edges only between occurrences of one metavariable, and a pattern occurrence $s$ of $M$ has $s\to x$ for every occurrence $x$ of $M$ (the pair is valid and $T^*$ covers $D$), which is (ii). $T_{\mathrm{all}}$ puts $M_K(Y_s)$ at $s$ and $M_K(\bar u_{s,x})$ at $x$, where by forcing and \evN\ $\bar u_{s,x}$ is the map $\bar y\mapsto\bar t_x$; so $T_{\mathrm{all}}$ is $T^*$ up to renaming and permutation of argument places. The construction of \Cref{prop:single:lgg} is polynomial by \Cref{cor:single:poly}.
\end{proof}

\noindent\emph{Computed.} The referee's brute force agreed with \Cref{prop:single:lgg} on all $115$ decided random cases, and in all $63$ anchor cases the lgg was the target.\src{referee t9}

% =====================================================================
\subsection{Computations}\label{app:single:computations}

All scripts are in \texttt{research/tracks/single/} and are run there with \texttt{python3 <script> [args]}; outputs are the \texttt{.out} files of the same name. The brute-force enumerator \texttt{dtenum.py} does not use the theory: it enumerates prefixes of the common prefix, fillers from a term pool and all sharings, and keeps the determinate covering templates. Agreement with it is therefore an independent check within its bounds. The referee's code in \texttt{referee\_checks/} is a second, fully independent implementation.

\paragraph{Main agreements.} $\Min(D)$ from saturated templates agreed with the brute-force enumerator in all $103$ completed random cases, and $\Feat(D)$ with the intersection of all enumerated covering templates on $7621$ queries; the referee's independent code, with bounds containing $T_0(D)$ and every $T_\varphi$, agreed on $125{,}109$ queries over $519$ data sets (\Cref{thm:single:min,thm:single:feat}). \Cref{thm:single:anchor} agreed with feature-based anchor status on $248/248$ random targets with one or two metavariables of all types (referee, independently: $540/540$, and $519/519$ against brute force), and \Cref{cor:single:special}(a), (c) on all of $571$ and about $2000$ random data sets.

\begin{sloppypar}
Scripts and what they check: \texttt{e1\_matching.py} (matching: unique matcher and recovery $1500/1500$, generality by freezing $400/400$, linear time up to $|s|=3.3\cdot10^6$); \texttt{e2\_examples.py} (\Cref{thm:single:min}, \Cref{ex:single:c82}, \Cref{prop:single:blowup}, \Cref{prop:single:weaker}, \Cref{cor:single:unions}(i)); \texttt{e3\_finitary.py 120 7} and \texttt{120 8 low} (\Cref{thm:single:min,thm:single:feat}, \Cref{prop:single:lgg} against brute force: $81/81$ and $22/22$ cases, $\Acc=\Feat$ on $5804$ and $1817$ queries); \texttt{e4\_anchor\_random.py 70 S 2}, $S=11,12,13$, and \texttt{60 14 2} (\Cref{thm:single:anchor}: prediction $=$ feature truth $248/248$; brute force decided $149$ and agreed on $144$, the other $5$ being artifacts of its bounds); \texttt{e5\_specializations.py} (\Cref{cor:single:special}: first-order $571/571$; induction $495$, Separation $495$, Replacement $498$, $\in$-induction $491$); \texttt{e6\_escalation.py 60 5} and \texttt{e6b\_escalation\_structured.py 1} (\Cref{thm:single:esc}, \Cref{conj:single:kill}); \texttt{e7\_rates.py} (circular, see above) and \texttt{e10\_rates\_revision.py} (\Cref{thm:single:rates}, \Cref{prop:single:kappa}, ground templates); \texttt{e8\_counterexamples.py} (\Cref{prop:single:weaker}; the chains of \Cref{sec:single:esc}); \texttt{e9\_union.py} and \texttt{e15\_ground\_union.py} (\Cref{prop:single:untagged-anchor,prop:single:union-verifier}); \texttt{e12\_c11\_refuted.py} (\Cref{prop:single:quadratic}); \texttt{e13\_escalation\_structure.py 3 40} (\Cref{prop:single:binderfree}, \Cref{lem:single:triangle}); \texttt{e16\_union\_lower\_bound.py} (\Cref{prop:single:union-esc}). The referee's scripts \texttt{t1}--\texttt{t10} check everything above except the items added in the revision, with full agreement ($\Feat=$ brute-force intersection on $125{,}109$ queries; anchor prediction $540/540$).
\end{sloppypar}

\noindent Scripts \texttt{e11} and \texttt{e14} check \Cref{thm:many:complexity} (track single, Theorem H).
